\section{Introduction}
\label{sec:intro}

When a human believes their work will be graded, they often think more carefully. This is the canonical social-facilitation effect~\citep{zajonc1965social}, and it is one of the few cognitive levers that depends only on \emph{framing}, not on the content of the task. Does the same lever exist for large language models?

The prompt-engineering literature gives a tantalising hint that it does. \citet{li2023emotionprompt} show that appending stimuli such as ``This is very important to my career'' lifts accuracy by roughly 10\% across 45 tasks. But the sycophancy literature pushes in the opposite direction: under a follow-up rebuttal such as ``Are you sure?'', frontier models flip roughly half of their answers and drop 17\% in accuracy~\citep{laban2023flipflop}, and they do so more strongly when the rebuttal is delivered \emph{before} the answer~\citep{fanous2025syceval}. These two threads point at opposite conclusions about the same intervention: a user signalling that they are paying close attention.

\para{The gap.} No prior work has tested a \emph{graded} preemptive skepticism manipulation that holds content fixed and varies only tone, measures both reasoning length and accuracy and capitulation, and ranges from neutral through curious, probing, judgmental, and hostile. The closest precedents each cover one slice: \citet{fanous2025syceval} vary rebuttal \emph{content strength} but not tone; \citet{kim2025challenging} show that casual phrasing sways models more than formal phrasing; \citet{hong2025sycon} find that a third-person ``Andrew prompt'' reduces sycophancy by 64\%. The hypothesis these works collectively suggest — that there exists a non-monotonic sweet spot where mild skepticism \emph{improves} reasoning without triggering capitulation — has not been tested directly.

\para{Our approach.} We construct a six-level tone ladder (Table~\ref{tab:tone_ladder}) as a sequence of system messages applied before identical user-message question content. We then run 3{,}900 single-turn calls and 1{,}200 follow-up rebuttal calls across two OpenAI models (\gptfo, \gptfomini), three datasets (\tqa, \gsm, \mmlu), and the full ladder. For every response we record accuracy, completion-token count, apology-keyword incidence (per the \psorry marker of \citet{laban2023flipflop}), and — for the two-turn experiment — the post-rebuttal flip outcome. The design lets us simultaneously test whether (i) tone modulates reasoning length, (ii) tone modulates accuracy non-monotonically, (iii) the operational apology marker the literature uses still functions, and (iv) preemptive tone primes resistance or fragility to a later rebuttal.

\para{What we find.} Three results stand out. First, \textbf{models do think longer when judged}. On \tqa, mean completion tokens rise by 36\% from neutral to ``I'm going to grade this. Don't be sycophantic'' on \gptfomini, and by 57\% on \gptfo (Wilcoxon paired $p \le 10^{-9}$). Second, the \textbf{single-turn ``sweet spot'' is real but small}: \gptfomini gains $+4$\,pp on \tqa under a pooled skeptic-majority vote (4 items recovered, 0 lost; one-sided exact $p=0.0625$); \gptfo peaks at $+4$\,pp at the same tone level; \gsm and \mmlu sit at ceiling and show no movement. Third, the \textbf{apology marker that the sycophancy literature defines as capitulation is extinct} in modern OpenAI models — exactly $0/5{,}100$ responses contain any apology keyword — but \textbf{preemptive judgmental tone makes the model more fragile} to a follow-up rebuttal: correct$\rightarrow$wrong flips on \tqa rise from 4\% at neutral to 13\% under ``I'm going to check your work carefully.''

\para{Contributions.}
\begin{itemize}[leftmargin=*,itemsep=0pt,topsep=0pt]
    \item We introduce a content-fixed six-level preemptive tone ladder and the first paired measurement of accuracy, completion-token expenditure, apology incidence, and rebuttal-induced flip rate as a function of tone.
    \item We provide the first direct evidence of an LLM analogue of social facilitation in reasoning length, with $p \le 10^{-9}$ for the strongest tone vs.\ neutral on \tqa across two frontier models.
    \item We document that the apology pattern \citet{laban2023flipflop} used to localise sycophantic capitulation is now silent in current OpenAI models, and that capitulation under rebuttal manifests instead as a confident silent rewrite — which has direct consequences for how the field should measure sycophancy going forward.
    \item We show that preemptive judgmental tone, far from priming downstream robustness, \emph{amplifies} rebuttal-induced fragility: the same tone region that boosts single-turn accuracy doubles or triples post-rebuttal correct$\rightarrow$wrong flips on \tqa.
\end{itemize}

\para{Paper organisation.} \secref{sec:related} positions the work against the sycophancy and prompt-tone literatures. \secref{sec:method} describes the ladder, datasets, models, and statistical plan. \secref{sec:results} reports per-cell accuracy, tokens, apology rate, and FlipFlop carry-over. \secref{sec:discussion} discusses where the sweet spot lives and why it is not a free lunch. \secref{sec:conclusion} summarises and outlines follow-up work.
